\section{总结与延伸}

本节把整期压缩为一张路线图。自动驾驶十年演进可以看成三次“升维”：第一次，从规则和高精地图走向 BEV/Transformer 的统一空间表征；第二次，从模块化功能开发走向端到端能力训练；第三次，从模仿人类驾驶轨迹走向 VLM、世界模型和 RL 的物理世界 Agent。

\begin{importantbox}{六个核心结论}
第一，高精地图和激光雷达路线解决的是短期确定性，但难以城市规模化。第二，Camera 信息密度高、成本低，是量产视觉路线的基础。第三，BEV 的关键是统一多相机空间表征，减少后融合错误。第四，端到端代表从功能开发到能力训练的范式转变。第五，L3/L4 的本质是责任边界，不是功能清单。第六，量产自动驾驶的长期壁垒来自真实数据、算力、组织执行和安全验证闭环。
\end{importantbox}

\subsection{关键术语速查表}

\noindent\begin{tabularx}{\textwidth}{p{0.22\textwidth}p{0.34\textwidth}X}
\toprule
术语 & 本期含义 & 课程关系 \\
\midrule
高精地图 & 预先记录道路几何和语义信息 & 早期确定性路线的核心依赖。 \\
LiDAR & 主动测距传感器 & 几何直接但成本和信息密度有约束。 \\
Camera & 被动视觉传感器 & 信息密度高，适合低成本量产和大规模数据。 \\
BEV & Bird's Eye View，鸟瞰空间表征 & 多相机统一建模的关键步骤。 \\
Transformer & 通用序列/注意力建模架构 & 被迁移到自动驾驶空间理解。 \\
端到端 & 从输入到轨迹/动作的模型化映射 & 自动驾驶从功能工程转向能力训练。 \\
VLM & Vision-Language Model，视觉语言模型 & 引入语义理解和复杂场景解释。 \\
World Model & 预测动作后果和未来世界状态的模型 & 支持规划、仿真和 RL 闭环。 \\
RL & Reinforcement Learning，强化学习 & 用反馈优化策略，是闭环自动驾驶方向之一。 \\
\bottomrule
\end{tabularx}

\subsection{后续观察问题}

\begin{enumerate}
\item 端到端路线能否在更多城市、天气、道路结构和长尾场景里稳定超过规则系统？
\item Camera-only、LiDAR-heavy 和多传感融合路线会如何在成本、安全和法规之间重新平衡？
\item 世界模型是否会成为自动驾驶下一阶段的核心，还是只作为训练/仿真辅助？
\item L3 在中国量产落地时，责任边界、接管机制和用户教育如何设计？
\item 主机厂的车队数据闭环，能否成为模型公司难以复制的长期壁垒？
\item 组织高压和安全验证如何平衡，避免为了上线速度牺牲可靠性？
\end{enumerate}

\subsection{拓展阅读}

\begin{itemize}
\item 对自动驾驶商业模式和 Waymo/Momenta 差异感兴趣，可对照 EP132 星海图高继扬访谈。
\item 对 Physical AI 和车端大模型工厂感兴趣，可对照 EP120 小鹏刘先明访谈。
\item 对 VLA、世界模型和机器人基础模型感兴趣，可对照 VLA 投屏版 `eiQFomOuCJs` 与 EP102 张祥雨多模态访谈。
\item 对硬件成本和激光雷达产业感兴趣，可对照 EP111 禾赛访谈。
\end{itemize}

\end{document}
